\scanpage{10}
toutes égales à $1$). On trouve alors une description complète de la
variation des valeurs propres mod $p$ de frobenius opérant sur $E(Y)$, pour
$Y$ variable avec un point de $U_d$, en termes essentiellement de la seule
représentation de $G$ dans un vectoriel de dimension finie sur $\underline F_p$, corres-
pondant aux $E(Y,\underline O_Y)_{\mathrm{SS}}$ (donc essentiellement aux ``matrices de Hasse-Witt\unsure{''}
des $Y$ en dimension critique), -- et bien entendu de la distribution des
\deleted{frobenius} éléments de frobenius dans $G$ en termes d'un point variable de
$U_d$.

8.5.10. L'homomorphisme fondamental $\phi$ de (8.5.1) est \deleted{surjectif}
\deleted{(resp.} injectif\deleted{)} \inserted{donc bijectif} si et seulement si $w$ dans 8.5.\unsure{1} est \deleted{injectif}
\deleted{(resp. si $w$ est} surjectif, \inserted{et, \unsure{il l'est si} $({}_p\operatorname{Tors}^n(X))_{\mathrm{ss}}=0$,} le seulement si étant vrai sous la réserve
$\operatorname{Ker}v=0$ de 8.5.7); il suffit en particulier pour ceci que $X$ n'ait pas
de $p$-torsion en \deleted{caractéristique} dimension $n\deleted{\unsure{-1}}$ \deleted{(resp. en dimension \unsure{$n$}\deleted{\unsure{ou}}
$n+1$)} ou ce qui revient au même par dualité, en dimension $n\unsure{+1}$). Lorsque $X$
se \deleted{Je pense que même si} relève (avec sa polarisation) en caractéristi-
que nulle, ce\deleted{s}\inserted{tte} condition\deleted{s} \deleted{\unsure{resp.}}\inserted{suffisante\deleted{s}} doi\deleted{vent}\inserted{t} aussi pouvoir s'exprimer (une
fois qu'on disposera de fondements satisfaisants pour la cohomologie
cristalline) par \unsure{l\deleted{es}\inserted{a}} condition\deleted{s} analogue\deleted{s} sur la \inserted{torsion de} cohomologie $p$-\uline{adique}
de la variété relevée (sur un corps algébriquement clos, ou si on veut,
sur le corps des complexes $\ldots$): \deleted{on réécrit (8.5.3) en termes de
$p$-torsion de cette variété tensorisée par $\underline Z/p\underline Z$ (mais en supprimant les
indices ss, qui n'ont plus de sens), et on demande que $w_p$ ainsi défini
soit surjectif resp. injectif $\ldots\unsure{.}$

Je pense que même si $n=\dim X=2$, il peut arriver que $\phi$ ne soit
pas surjectif, \deleted{i.e.} qu'on ne peut pas dans l'inégalité (8.5.6) supprimer
le terme correctif $\dim\operatorname{Ker}w$ ou $t^{n-1}(X,p)$. En effet\deleted{\unsure{,}}, si $X$ est une}
\mspage{9}
